% DOCUMENT FORMAT ======================= -*- mode: LaTeX; coding: utf-8 -*- ===

\documentclass[diploma]{softlab-thesis}


% PACKAGE SETTINGS =============================================================

\usepackage[cm-default]{fontspec}
\usepackage{amsmath}
\usepackage{amsfonts}
\usepackage{multirow}
\usepackage{array}
\usepackage{mdwlist}
\usepackage{subfig}
\usepackage{floatrow}
%\usepackage{float}
\usepackage{verbatim}
\usepackage{color}
\usepackage{graphicx}
\usepackage{xunicode}
\usepackage{xltxtra}
\usepackage{url}
%\usepackage{dsfont}
%\usepackage{microtype}
\usepackage[htt]{hyphenat}
\usepackage{multicol}
\usepackage{wrapfig}

% FONT SETTINGS ===============================================================

%\setromanfont[Mapping=tex-text]{CMU Serif}
%%\setromanfont[Mapping=tex-text]{CMU Sans Serif} % temporary change until printing
%%\setsansfont[Mapping=tex-text]{CMU Sans Serif}
%%\setmonofont[Mapping=tex-text]{CMU Typewriter Text}
%\setmainfont[Mapping=tex-text]{CMU Serif}
%%\setmainfont[Mapping=tex-text]{CMU Sans Serif}  % temporary change until printing

%\setromanfont[Mapping=tex-text,ExternalLocation=fonts/]{cmunrm.otf}
%\setsansfont[Mapping=tex-text,ExternalLocation=fonts/]{cmunss.otf}
%\setmonofont[Mapping=tex-text,ExternalLocation=fonts/]{cmuntt.otf}
%\setmainfont[Mapping=tex-text, ExternalLocation=fonts/]{cmunss.otf}

% CUSTOM COMMANDS ==============================================================

% Layout macros
\newcommand\spa[1]{\; #1 \;}

% Font macros
\newcommand\resfont[1]{\ensuremath{\mathtt{#1}}}

% Mathematical macros
\newcommand\setmap[3]{#1\{#2 \mapsto #3\}}
\newcommand\getmap[3]{(#2 \mapsto #3) \in #1}
\newcommand\tuple[2]{\ensuremath{\langle#1, #2\rangle}}
\newcommand\mfrac[2]{\ensuremath{\dfrac{#1}{#2}}}
\newcommand\nequiv[2]{\ensuremath{#1 \not\equiv #2}}

% Core Ruby Operational Semantics letter bindings
\newcommand\mem{\mu}

% Core Ruby Operational Semantics low level macros
\newcommand\state[2]{(#1, #2)}
\newcommand\transition[1]{\ensuremath{\overset{#1, c*}{\rightarrow}}}
\newcommand\range[2]{#1, ..., #2}
\newcommand\midrange[5]{\range{#1}{#2}, #3, \range{#4}{#5}}

% Core Ruby Operational Semantics high level macros
\newcommand\operation[5]{\ensuremath{\state{#1}{#2} \transition{#3} \state{#4}{#5}}}
\newcommand\propagation[2]{\operation{#1}{\mem}{#2}{#1'}{\mem'}}

% Core Ruby specific Operational Semantics macros
\newcommand\semicolon[2]{#1; \; #2}
\newcommand\assign[2]{#1 = #2}
\newcommand\mcall[3]{#1.\texttt{#2}(#3)}
\newcommand\ifte[3]{\resfont{if} \; #1 \; \resfont{then} \; #2 \; \resfont{else} \; #3}
\newcommand\newclass[2]{#1.\resfont{new}(#2)}
\newcommand\methoddef[3]{\resfont{def} \; #1(#2) = #3}
\newcommand\classdef[2]{\resfont{class} \; #1 = #2}
\newcommand\with[3]{with \; \tuple{#1}{#2} \; do \; #3}

% Success Typing macros
\newcommand\ssub{\sqsubseteq_S}

% Core Ruby Success Typing letter bindings
\newcommand\classlist{\Delta}
\newcommand\envir{\Gamma}
\newcommand\fields{\Phi}
\newcommand\currclass{l}

% Core Ruby Success Typing inferencing macros
\newcommand\stinfer[5]{\classlist; \; #1; \; \fields \; \underset{\currclass}{\vdash} \; #2: #3 \; \& \; #4; \; #5}


%%%%%%%%%%%%%%%%%%%%%%%%%%% HASKELL STUFF %%%%%%%%%%%%%%%%%%%%%%%%%%%

%\newcommand\typerep[1]{\ensuremath{#1}}
\newcommand\typerep[1]{\lstinline[basicstyle=\normalsize\ttfamily,keywords={}]|#1|}
\newcommand\typefootrep[1]{\textbf{\lstinline[basicstyle=\footnotesize\ttfamily,keywords={}]|#1|}}
% \newcommand\ttyperep[1]{\typerep{#1}}
% \newcommand\mtyperep[1]{\mbox{\typerep{#1}}}

% Arrow types
\newcommand\typeto[2]{\typerep{#1} \typerep{->} \typerep{#2}}
\newcommand\typetotwo[3]{\ensuremath{\typerep{#1} \typerep{->}
                                     \typerep{#2} \typerep{->}
                                     \typerep{#3}}}

\newcommand\tyconapone[2]{\ensuremath{\mbox{\typerep{#1}} \:\: \mbox{#2}}}
\newcommand\tyconaponeC[2]{\ensuremath{\mbox{\typerep{#1}} \:\: \mbox{\typerep{#2}}}}
% \newcommand\tyconapone[2]{\typerep{#1} $\:$ \typerep{#2}}

\newcommand\tyconaptwo[3]{\ensuremath{\mbox{\typerep{#1}} \:\: \mbox{#2} \:\: \mbox{#3}}}


% FLOAT SETUP ==================================================================

% SPELLING =====================================================================

% CODE HIGHLIGHTING ============================================================

\usepackage{listings}

\definecolor{gray}{rgb}{0.5,0.5,0.5}
\definecolor{darkgreen}{rgb}{0.0,0.5,0.0}

\lstset{
	language=Haskell,
	basicstyle=\footnotesize\ttfamily,
	numbers=left,
	numberstyle=\tiny\color{gray},
	stepnumber=1,
	numbersep=5pt,
	backgroundcolor=\color{white},
	showspaces=false,
	showstringspaces=false,
	showtabs=false,
	frame=single,
	rulecolor=\color{black},
	tabsize=2,
	captionpos=b,
	breaklines=true,
	breakatwhitespace=false,
	title=\lstname,
        xleftmargin=17pt,
        xrightmargin=5pt,
        aboveskip=\baselineskip
}

\newcommand\haskellcode[2]{\lstinputlisting[language=Haskell,caption={#1},label=#1]{src/#2}}
\newcommand\rubycode[2]{\lstinputlisting[language=Ruby,caption={#1},label=#1]{src/#2}}
\newcommand\erlangcode[1]{\lstinputlisting[language=erlang,captionpos=]{src/#1}}
\newcommand\plaintext[1]{\lstinputlisting[language=,captionpos=,stepnumber=0]{src/#1}}

% CHANGE MATH FONT ============================================================

% HYPERREF MUST BE LAST =======================================================

\usepackage[xetex,colorlinks=true,linkcolor=blue,citecolor=darkgreen]{hyperref}

% DOCUMENT INFORMATION =========================================================

\title{Thesis subject} % ===============> FIXME
\author{Όνομα Φοιτητή}
\authoren{Student name}
\datedefense{9}{9}{9999}      % ===============> FIXME
\TRnumber{CSD-SW-TR-*-*} % ok ==================> FIXME
\supervisor{Υπέυθυνος Διπλωματικής}
\supervisorpos{Τίτλος Υπευθύνου}
\committeeone{Πρώτο μέλος επιτροπής}
\committeeonepos{Τίτλος μέλους}
\committeetwo{Δεύτερο μέλος επιτροπής}
\committeetwopos{Τίτλος μέλους}
\committeethree{Τρίτο μέλος επιτροπής}
\committeethreepos{Τίτλος μέλους}

\hypersetup{
	pdftitle={Thesis subject}, % ===> FIXME
	pdfauthor={Όνομα Φοιτητή},
	pdfsubject={},
	pdfkeywords={}
}


% MAIN DOCUMENT ================================================================

\begin{document}

\frontmatter
\maketitle

\def\templen{\parindent}
\setlength{\parindent}{0pt}
\setlength{\parskip}{1.5ex plus 0.5ex minus 0.2ex}
% DOCUMENT FORMAT ======================= -*- mode: LaTeX; coding: utf-8 -*- ===
\begin{abstractgr}
	Περίληψη της διπλωματικής.
\begin{keywordsgr}
  	Λέξη-κλειδί 1, λέξη-κλειδί 2, λέξη-κλειδί 3
\end{keywordsgr}
\end{abstractgr}

\begin{abstracten}
	Abstract of diploma thesis.
\begin{keywordsen}
  	Key-word 1, Key-word 2, Key-word 3
\end{keywordsen}
\end{abstracten}

\begin{acknowledgementsgr}
	Ευχαριστίες.
\end{acknowledgementsgr}





\selectlanguage{english}   % Everything from here on is in english
\setlanguage{english}
\nogreekalph

\setlength{\parindent}{\templen}
\setlength{\parskip}{0pt}
\tableofcontents
\listoffigures
%\listoftables
\renewcommand{\lstlistlistingname}{List of Listings}
%\lstlistoflistings % changed the title above

\mainmatter
% moved these two commands here so that they don't influence the toc
\setlength{\parindent}{0pt}
\setlength{\parskip}{1.5ex plus 0.5ex minus 0.2ex}

\renewcommand\floatpagefraction{.7}

\chapter{Introduction}\label{ch:intro}

\section{Introduction/Motivation}

Bla-bla...

\section{Thesis structure}

% TODO: Make `Chapter x` click-able and anchored to the respective chapter
\begin{description}
\item[Chapter~\ref{ch:ch2}:]
We define what "cloud" means and mention some of the most notable examples.
Then, we give a brief overview of the synnefo implementation, its key
characteristics and why it can have a place in the current cloud world.
\item[Chapter~\ref{ch:ch3}:]
We present the architecture of Archipelago and provide the necessary theoretical
background (mmap, IPC) the reader needs to understand its basic concepts. Then,
we thoroughly explain how Archipelago handles I/O requests. Finally, we mention
what are the current storage mechanisms for Archipelago and evaluate their
performance.
\item[Chapter~\ref{ch:ch4}:]
We explain why tiering is important and what is the state of tiered storage at
the moment (bcache, flashcache, memcached, ramcloud, couchbase).  Then, we
provide the related theoretical background for cached (hash-tables, LRUs).
Finally, we defend why we chose to roll out our own implementation.
\item[Chapter~\ref{ch:cached-design}:]
We explain the design of cached, the building blocks that is consisted of
(xcache, xworkq, xwaitq). Then, we provide extensive benchmark results and
compare them to the ones of Chapter 3.
\item[Chapter~\ref{ch:ch6}:]
TODO
\item[Chapter~\ref{ch:ch7}:]
We draw some concluding remarks and propose some future work.
\end{description}

\chapter{Chapter 2}\label{ch:ch2}

\section{Section 1}

This section has an important citation\cite{inproceeding}

\subsection{Subsection 1}

This subsection has code in Haskell:
\haskellcode{Sample code}{samplecode.hs}

It also has a list:
\begin{description}
\item[Item 1]
	First item
\item[Item 2]
	Second item and a footnote\footnote[1]
	{Footnote description.}.
\item[Item 3]
	Third item and text in \textit{italics}.
\end{description}

And an enumerated list:

\begin{enumerate}
\item	First item.
\item	Second item and text in \textbf{bold}
\end{enumerate}

\section{Section 2}

\subsection{Subsection 1}
This subsection has a link to the block of code~\ref{Sample code} in Section 1.

\subsection{Subsection 2}

% FIXME: Must add content here
This subsection has a FIXME comment, visible only to the author.

\chapter{Chapter 3}\label{ch:ch3}

\section{Section 1}

This how we add a url: \url{http://www.example.org}

\section{Section 2}

And this is how we point to Figure~\ref{fig:image1}.

\begin{figure}[hb]
	\centering
	\includegraphics[]{diagrams/9.jpg}
	\caption{This is an image}
	\label{fig:image1}
\end{figure}

\subsection{Sub-section 3}

Another way to create a list:
\begin{itemize}
	\item Item 1
	\item Item 2
	\item Item 3
	\item Item 4
\end{itemize}


\chapter{Chapter 4}\label{ch:ch4}

\section{Section 1}

We can also use a special fonts to differentiate between text and $math()$.


\chapter{Design of cached}\label{ch:cached-design}

\section{General}

Currently, Archipelago doesn't have an xseg peer that can act as a cache for
block requests. Having a cache however is essential for a tiered storage and
cached peer is implemented for this reason.

\section{Implementation details}

cached is based on xseg and also on the following xtypes:

\begin{itemize}
	\item \textbf{xcache} (for cache support)
	* xwork (for job support)
	* xworkq (for atomicity in execution of jobs)
	* xwaitq (for conditional execution of jobs)
\end{itemize}

More specifically, cached consists of the cache provided by xcache and a
pre-allocated number of objects. An object is divided in buckets and its size,
as well as bucket size, are defined by the user.

The fact that objects are pre-allocated means two things:

1) We don't need to care about memory fragmentation and system call overhead
2) We cannot index single buckets. <FILLME>

\section{The xcache xtype}



\section{The xworkq xtype}

Every object has a workq. Whenever a new request is accepted/received for an
object, it is enqueued in the workq and we are sure that only one thread at a
time can have access to the objects data and metadata.

For more information, see the xworkq.

\section{The xwaitq xtype}

When a thread tries to insert an object in cache but fails, due to the fact
that cache is full, the request is enqueued in the xcache waitq, which is
signaled every time an object is freed.

For more information, see the xwaitq.

\section{Cached internals}

\subsection{Object states}

Every object has a state, which is set atomically by threads. The state list is
the following:

* READY: the object is ready to be used
* FLUSHING: the object is flushing its dirty buckets
* DELETING: there is a delete request that has been sent to the blocker for
  this object
* INVALIDATED: the object has been deleted
* FAILED: something went very wrong with this object

Also, object buckets have their own states too:

* INVALID: the same as empty
* LOADING: there is a pending read to blocker for this bucket
* VALID: the bucket is clean and can be read
* DIRTY: the bucket can be read but its contents have not been written to the
  underlying storage
* WRITING: there is a pending write to blocker for this bucket

Finally, for every object there are bucket state counters, which are increased/
decreased when a bucket state is changed. These counters give us an O(1)
glimpse to the bucket states of an object.

\subsection{Per-object peer requests}

Reads and writes to objects are practically read/write request from other
peers, for which a peer request has been allocated. There are cases though
when an object has to allocate its own peer request e.g. due to a flushing of
its dirty buckets. Since this must be fast, there are pre-allocated requests
hard-coded in the struct of each object which can be used in such cases.

\subsection{Write policy}

The user must define beforehand what is the write policy of cache. There are
two options: writethrough and writeback. On a side note, as far as reads and
cache misses are concerned, cached operates under a write-allocate policy.



\backmatter
\cleardoublepage % start at the next odd page
\phantomsection % correct hyperlinking
\bibliography{references}
\bibliographystyle{softlab-thesis}
% \include{glossary}
% \chapter{Appendix}
% \printindex

\end{document}
